% eksperymentalne tłumaczenie na włoski

%

\pol{\section{Problem Snelliusa-Pothenota}}
\ita{\section{Problema di Snellius-Pothenot}}

\pol{Problem Snelliusa-Pothenota rozwiąże po raz pierwszy Willebrord Snellius w 1615 roku.}
\ita{Il problema di Snellius-Pothenot sarà risolto per la prima volta da Willebrord Snellius nel 1615.}
\pol{\index{problem!Snelliusa-Pothenota}}%
\ita{\index{problema!di Snellius-Pothenot}}%

\pol{Wkład Laurenta Pothenota (żyjącego w latach 1650--1732) w problem lub jego rozwiązanie jest niejasny, dlatego George McCaw, słynny brytyjski geodeta stwierdzi, że powinniśmy pomijać jego nazwisko.}
\ita{Il contributo di Laurent Pothenot (1650--1732) al problema o alla sua soluzione non è chiaro, perciò George McCaw, celebre geodeta britannico, sosterrà che dovremmo omettere il suo nome.}

\begin{problem}
    \pol{Dane są trzy znane punkty $A$, $B$ i $C$ oraz nieznany punkt $P$.}
    \ita{Sono dati tre punti noti $A$, $B$ e $C$ e un punto ignoto $P$.}
    \pol{Z punktu $P$ obserwator mierzy kąty utworzone przez linie łączące go z każdym z pozostałych trzech punktów.}
    \ita{Dal punto $P$ un osservatore misura gli angoli formati dalle linee che lo congiungono con ciascuno degli altri tre punti.}
    \pol{Wyznaczyć położenie punktów $P$.}
    \ita{Determinare la posizione del punto $P$.}
\end{problem}

\pol{Zdegenerowany przypadek, kiedy punkt $P$ znajduje się blisko okręgu opisanego na trójkącie $ABC$ powinien być unikany jak ogień, szczególnie jeśli $P$ oznacza położenie rozbitka na morzu, zaś $A$, $B$, $C$ to trzy latarnie na lądzie.}
\ita{Il caso degenere in cui il punto $P$ si trovi vicino al cerchio circoscritto al triangolo $ABC$ dovrebbe essere evitato come la peste, soprattutto se $P$ rappresenta la posizione di un naufrago in mare e $A$, $B$, $C$ sono tre fari sulla costa.}
%
